\documentclass[10pt,a4paper]{amsart}
\usepackage[margin=19mm]{geometry}
\usepackage{amsmath,amssymb,mathtools}
\usepackage[scr=rsfs,cal=euler]{mathalfa}
\usepackage{xfrac}
\usepackage[unicode,hidelinks,bookmarks=false]{hyperref}
\newcommand{\ie}{i.e.\ }
\newcommand{\cf}{cf.\ }
\newcommand{\resp}{respectively\ }
\newcommand{\Subsection}{\subsection}
\providecommand{\nobreakdash}{\nobreak\hbox{-}}
\setlength{\emergencystretch}{2em}
\multlinegap0pt
\usepackage{amssymb}
\usepackage{esint}
\usepackage[all]{xy}
\usepackage{xcolor}
\newtheorem{lem}{Lemma}
\newtheorem{prop}{Proposition}
\def\Area{\mathrm{Area}}
\title{Stability for the 3D Riemannian Penrose inequality}
\author{Conghan Dong}
\date{Source: arXiv:2402.10299v4 (CC BY 4.0)}
\begin{document}
\maketitle

\noindent\emph{Source and licence.} Stability for the 3D Riemannian Penrose inequality. Version arXiv:2402.10299v4 (CC BY 4.0), distributed by arXiv under the Creative Commons Attribution 4.0 International licence, http://creativecommons.org/licenses/by/4.0/, as recorded in the arXiv licence metadata for this exact version and shown on the abstract page https://arxiv.org/abs/2402.10299v4. Full licence notice: \texttt{licenses/arxiv-2402.10299-CC-BY-4.0.txt}.\par\smallskip

\noindent\textit{Editorial scope.} Counted unit: the article's prop f2-eq (source0.tex lines 1443--1448). The article's complete original proof of it is delivered with it (source0.tex lines 1449--1475, one \texttt{proof} environment of 27 lines). No mathematics is written, repaired or re-derived: the statement and the proof are the article's own lines, in the article's own notation, and no retained line is substituted or re-flowed.  Retained uncounted as the source-local context the proof needs: lines 1397--1399; lines 1417--1423.  Nothing was omitted from any retained span, so the package carries no omission record.  No retained line contains a citation, so the wrapper carries no bibliography and no bibliography entry.  No retained line contains a figure environment, an image-inclusion command or drawing code, so no image is delivered and the package declares no asset dependency.  Editorial formatting only, in addition: an AMS \texttt{amsart} preamble that restates the article's own declarations and macros used by the retained lines, namely the environments \texttt{lem}, \texttt{prop} on separate counters, together with the standard packages those lines need.  The retained environments print in this document as Lemma 1, Proposition 1 in order of appearance, whereas the article prints the same items with the numbering of its own full document.  The complete native source of the article as distributed in the arXiv e-print for this exact version is bundled unchanged as \texttt{sources/154644/source0.tex}.
\begin{align}\label{grad-f}
\left| |\nabla f|_h -  \frac{2(1-f)^2}{\mathcal{C}(\Sigma ,g)} \right| \leq \Psi (m(h) ).	
\end{align}

\begin{lem}\label{weak-W12}
	Up to a subsequence, $f_i$ converges to $f_\infty$ in the weakly $W^{1,2}$-sense. That is, for any uniformly converging sequence of compactly supported Lipschitz functions $\psi _i \to \psi $ and $\nabla \psi _i \to \nabla \psi $ in $L^2$, we have 
	\begin{align*}
		\lim_{i\to \infty} \int_{\mathcal{E}_i} f _i \psi _i \mathrm{dvol}_{h_i} & = \int_{\mathbb{R}^3} f _\infty \psi \mathrm{dvol}_{\mathrm{Eucl}},\\
		\lim_{i\to \infty} \int_{\mathcal{E}_i} \left<\nabla f _i, \nabla \psi _i \right>_{h_i} \mathrm{dvol}_{h_i} &= \int_{\mathbb{R}^3} \left<\nabla f _\infty, \nabla \psi  \right>_{\mathrm{Eucl} } \mathrm{dvol}_{\mathrm{Eucl}}.
	\end{align*}
\end{lem}

\begin{prop}
	$f _\infty$ satisfies the following equation weakly on $\mathbb{R}^3$ :
	\begin{align}\label{f2-eq}
\Delta _{\mathrm{Eucl} } f_\infty ^2 = \frac{24}{m_0^2} \cdot (1-f_\infty)^4.
\end{align}
\end{prop}

\begin{proof}
It's enough to show that for any smooth function $\psi $ on $\mathbb{R}^3$ with compact support in $U \subset B(x_o, D) $, 
\begin{align*}
	-\int_{U} \left<\nabla f^2 _\infty, \nabla \psi\right>_{\mathrm{Eucl} } \mathrm{dvol}_{\mathrm{Eucl}} = \frac{24}{m_0^2} \int_{U} (1-f_\infty)^4 \psi \mathrm{dvol}_{\mathrm{Eucl}} .
\end{align*}
Define $\psi _i := \psi \circ \mathbf{u}_i$ on $\mathcal{E}_i$. Then $\psi _i$ are uniformly Lipschitz and $\psi _i \to \psi $ uniformly.

Notice that
\begin{align*}
	\Delta _{h_i}f_i^2 = 6|\nabla f_i|_{h_i}^2.
\end{align*}
Using integration by parts, for $U_i := \mathbf{u}_i^{-1}U \subset \mathcal{E}_i$, we have
\begin{align*}
	-\int_{U_i} \left< \nabla f _i^2, \nabla \psi _i \right>_{h_i} \mathrm{dvol}_{h_i} &= \int_{U_i} 6 \psi _i |\nabla f_i|^2_{h_i} \mathrm{dvol}_{h_i}  - \int_{\partial U_i} \psi _i \left< \nabla f _i, \vec{n} \right> \mathrm{dA}_{h_i}.
\end{align*}
Since $|\psi _i |, |\nabla f _i|_{h_i} \leq C$ on $U_i$,
$$
\left|\int_{\partial U_i} \psi_i \left<\nabla f _i, \vec{n} \right> \mathrm{dA}_{h_i}\right|\leq C\cdot  \Area_{h_i}(\partial \mathcal{E}_i) \to 0.
$$ 
Using (\ref{grad-f}), we have 
\begin{align*}
	\lim_{i\to \infty} \int_{U_i} 6 \psi _i |\nabla f_i|_{h_i}^2 \mathrm{dvol}_{h_i} &= \lim_{i \to \infty} \int_{U_i} \frac{24}{\mathcal{C}(\Sigma _i, g_i)^2} (1-f_i)^{4} \psi _i \\										  &= \frac{24}{m_0^2} \int_{U}(1-f_\infty)^4 \psi \mathrm{dvol}_{\mathrm{Eucl}} .
\end{align*}
From Lemma \ref{weak-W12}, we know 
$$\lim_{i\to \infty} \int_{U_i} \left< \nabla f^2 _i, \nabla \psi _i\right>_{h_i} \mathrm{dvol}_{h_i} = \int_{U} \left< \nabla f^2 _\infty, \nabla \psi\right>_{\mathrm{Eucl}} \mathrm{dvol}_{\mathrm{Eucl}} ,$$
which concludes the proof. 
\end{proof}

\end{document}
